\newcommand{\makefrontpage}
{\begin{titlepage}%
    \begin{pspicture}(0cm,0cm)(\textwidth,\textheight)

      \rput[tr](14.4,21.8){\includegraphics[scale=1]{sdulogo.png}}

      \rput[tl](0.01\textwidth,0.90\textheight){
        \begin{minipage}[t]{14cm}
         \textsf{\Large{\textbf{}}}\\
          \vspace*{0.10\textheight}
          \textit{\textsf{\Large{
                \begingroup
                \tabcolsep -2mm
                \begin{tabular}{ p{2,4cm} l }
   \textsf{\Large{\textbf{DAT2}}}             & \textsf{\Large{\textbf{Gruppe 2}}}\\
   &   \\ 
260387 & Rudi Hansen \\
050379 & Niels H�gh \\
280786 & Kelvin Kj�rvik Pagels \\
160788 & Kim Lindberg Schwaner\\
                \end{tabular}
                \endgroup}}}\\
          \begin{flushleft}
            \Huge{\textsf{\textbf{Pole Position}}}
          \end{flushleft}
          
          \vspace*{0.01\textheight}
          
          \begin{flushleft}
            \LARGE{\textsf{\textbf{Udvikling af system til at drive en Scalextricbil autonomt ved hj�lp af diverse sensorinput}}}
          \end{flushleft}

          \vspace*{0.04\textheight}
          \textsf{\Large{februar - juni 2011}}
          \vspace*{0.04\textheight}

          \textsf{\Large{Vejleder: Ren� Lynge Eriksen}}
           
          \vspace*{0.06\textheight}
              \textsf{I denne 2. semesters rapport unders�ges et design af styringen til en autonom scalextric bil. I rapporten gennemg�s hvilke kriterier sensorer og andre elektriske komponenter er blevet udvalgt efter. Deres virkem�de og hvordan de anvendes beskrives. Den udviklede software til ATmega32 microcontrolleren, til at styre bilen, gennemg�s ligeledes. Slutteligt konkluderes der at det, p� grund af uventede komplikationer med sensor input, ikke har v�ret muligt at f� bilen til at k�re autonomt.}\\
          
          

	\vspace*{0.05\textheight}
              \textsf{Det Tekniske Fakultet\\
                  Syddansk Universitet\\
                  Niels Bohrs All� 1\\
                  DK-5230 Odense M\\
                  \vspace*{0.3cm}Danmark\\
                  www.sdu.dk/tek\\
                  Tlf: (+45) 65 50 73 03\\
                  E-mail: tek@tek.sdu.dk\\
              }

        \end{minipage}}
    \end{pspicture}
   \end{titlepage}%
  \setcounter{footnote}{0}
  % 
 }
